% 第 38 章骨架：导引框已定，正文待写。写作计划见 docs/PLAN.md“第 38 章”，写法见 docs/STYLE.md。
\chapter{菲利普斯曲线与宏观计量模型}\label{ch:phillips}

\begin{chapterguide}
\guidequestion{通胀与失业真是一张可以点菜的菜单吗？}
\begin{guidegoals}
\item 复现菲利普斯（1958）的工资—失业关系，并说明他的数据处理方法；
\item 解释萨缪尔森与索洛（1960）怎样把它变成政策菜单；
\item 描述克莱因—戈德伯格模型（1955）等大型宏观计量模型的结构；
\item 指出这套框架隐含的假设，为第七部分做准备。
\end{guidegoals}
\guideprereq{\chref{ch:islm}\enspace\nameref{ch:islm}}
\end{chapterguide}

\chapterstub
